\section{Provenance, Measurement, and Intervention Audit}
\label{app:audit}
\subsection{Population and missingness}
\begin{table*}[t]
\centering\small\begin{tabular}{lrrrrr}
\toprule Model & Completed & Usable & Included (\%) & Unfinished & Segments (\%) \\
\midrule
DeepSeek-R1-Distill-Llama-8B & 3,281 & 3,297 & 94.5 & 18 & 100.00 \\
Gemma4-E4B-it & 3,488 & 3,217 & 92.2 & 0 & 99.99 \\
Gemma4-12B-it & 3,181 & 3,179 & 91.1 & 1 & 99.94 \\
Gemma4-31B-it & 3,484 & 3,483 & 99.9 & 0 & 99.99 \\
Qwen3.5-4B & 2,333 & 2,347 & 67.3 & 14 & 99.98 \\
Qwen3.5-9B & 2,915 & 2,925 & 83.9 & 10 & 99.99 \\
Qwen3.5-27B & 3,352 & 3,355 & 96.2 & 3 & 99.99 \\
\bottomrule
\end{tabular}

\caption{Atlas inclusion, out of 3,488 attempted generations per reasoning model. ``Usable'' means at least one segment passes both contextual and isolated parsing checks. ``Unfinished'' counts unfinished attempts within the usable population. The last column is the fraction of annotated segments retained, not coverage of unannotated text. The companion \protect\path{coverage.csv} reports all 56 model--domain denominators, including attempted, completed, annotated, usable, and excluded counts.}
\label{tab:coverage}
\end{table*}
Completion and annotation inclusion are distinct. The 24,416 reasoning attempts yield 21,803 usable traces, of which 46 are unfinished. Table~\ref{tab:coverage} reports the model totals. Missing annotation rows and entirely invalid annotations are counted separately; absence from the annotation table does not identify whether a failure occurred during generation, staging, or judging. High retention among annotated segments cannot establish representativeness of traces that were never annotated. The fixed-prefix multi-judge protocol below explicitly samples from attempted generations.

\subsection{Fitting, selection, and evaluation populations}
\begin{table*}[t]
\centering\small
\begin{tabularx}{\textwidth}{@{}lXX@{}}
\toprule Component & Verified source or procedure & Independence and scope\\\midrule
Behavioral readout & 800 prompts, 3,200 traces, 251,420 segment states; four domains, 200 prompts and four generations each. Stored domain means reproduced exactly. & Math, GPQA, and planning fitting prompts have zero ID or normalized-text overlap with their evaluation sets.\\
Domain contrasts & Five 16-target means; moral corpus added for profile estimation. Equal-domain reference mean; same saved map. & Reference estimation and readout fitting are separate steps; no target-specific dose calibration.\\
A2 archetype & Four medoids per domain; math pool filtered by historical correctness. Saved pilots compare A2 and A4 at multiple gains, with random controls. & Two 40-item pilots exclude final math items. Later comparisons reuse the 199-item evaluation set.\\
Layer, rank, dose & Layer 16, rank eight, ridge $10^5$; open-loop norm 6.75 after readout comparisons and causal pilots. & Historically selected settings; new validation is retrospective, not nested selection.\\
Historical reward-trained biases & 301 training prompts; one run per initialization. Saved configurations, metrics, and vectors retained. & No overlap with 199 evaluation items; training includes readout and SEAL-validation prompts. Per-rollout rewards and per-run source hashes absent.\\
Final evaluations & 199 math, 200 GPQA, 200 planning items; two decoding seeds, shared baselines verified by run fingerprint. & Holdout from specified fitting/pilot sets, but not a fresh test after every intervention-design decision.\\
\bottomrule
\end{tabularx}
\caption{Verified provenance and the resulting inferential boundaries. Membership is checked against saved task files; chronological notes are supporting records, not an immutable preregistration.}
\label{tab:provenance}
\end{table*}
The historical readout is identified by matching its stored domain means to all eight fitting shards with maximum absolute discrepancy zero. Input and map hashes accompany the audit. Prompt IDs and whitespace-normalized text establish the exclusions in Table~\ref{tab:provenance}. The 199 math problems were already used to compare open-loop and closed-loop A2 before the longer-budget study, so that study is not a fresh replication of intervention selection. SEAL layer selection uses a separate 101-item validation set. These distinctions prevent the 301/199 RL split from being mistaken for independence from every earlier decision.

\paragraph{What A2 denotes.}
Each corpus trace is represented by 24 rates: eight overall, eight first-half, and eight second-half rates. The saved construction requires completion, no hard-loop flag, and at least 20 valid segments; math additionally requires success under the historical grader. Coordinates are standardized within domain, and four Euclidean PAM medoids are formed, subsampling to at most 600 candidates with seed zero. A medoid's eight overall rates are repeated into the future and next-four target blocks. A2 is one selected math medoid, characterized by more subgoals and less verification/self-questioning relative to the math reference. The saved first pilot tests A2 and A4 at gains one and two; the second tests open and closed control at gains two and three. A2's later selection is outcome-informed. Regrading evaluation does not determine which medoid a newly graded construction would select.

\subsection{Readout validation}
The historical fitting source partitions by trace ID, allowing repeated generations of the same prompt in different folds, and estimates activation normalization and domain target means before partitioning. Those scores are not used as clean held-out validation. The revised analysis groups all generations and positions of a prompt, fits activation normalization and domain reference means on training data only, and fixes the historical layer, rank, and ridge. Five prompt folds test within-domain prediction. Four leave-domain-out fits use the training grand mean for an unseen domain; held-out target labels never define its reference. Mean squared errors first average positions within trace and traces within prompt. The reported skill is $1-\mathrm{MSE}_{\rm readout}/\mathrm{MSE}_{\rm reference}$, with prompt-bootstrap intervals. The saved intervention map is unchanged. This validates prediction separately from the manipulation checks below.
\IfFileExists{tables/readout_result.tex}{Across the 16 targets, mean prompt-balanced MSE reduction over the training-domain reference is 7.61\%. With an unseen domain, mean reduction over the training-grand-mean reference is 3.43\%; these baselines differ. The mean pooled position-level $R^2$ across targets is -0.068 for held-out domains, indicating poor absolute out-of-domain prediction despite improvement on some components. Figure~\ref{fig:readout-validation} gives per-target prompt-bootstrap intervals. Hyperparameters were selected historically, so these fixed-setting checks are not nested estimates of a model-selection procedure.
\begin{figure*}[t]
\centering\includegraphics[width=.95\textwidth]{figures/figA7_readout_validation.pdf}
\caption{Readout prediction under prompt-disjoint and held-domain validation. Skill compares MSE with the training-domain mean for known domains and the training grand mean for an unseen domain. Whiskers resample prompts; targets are future or next-four-segment rates. The historical steering map remains frozen.}
\label{fig:readout-validation}
\end{figure*}
}{\draftnote{Prompt-disjoint readout validation is running; no new score is claimed here.}}

\subsection{Grading and historical RL implementation}
The earlier deterministic grader could misread formatting or match fragments of mathematical expressions. Corrected evaluation uses reference-blind free-text extraction, checks that the extracted answer is grounded in the final response, and applies formatting-normalized mathematical equivalence. Unfinished attempts score zero. Reconciling the two saved step-60 evaluations changes 47 cutoff rows for the zero-start vector and 46 for the A2-initialized vector: nine traces and five unique problems in each run. These are evaluation corrections, not counts of affected training rewards. The saved aggregate training logs cannot recover the latter.

The inspected training source binds a log-probability closure to a variable named \texttt{T}, initially the sampling temperature, then reassigns \texttt{T} to sequence length before calling that closure. This is a source-level defect, although the absence of per-run source hashes prevents proving the precise executed version. Consequently, the paper retains the saved vectors as fixed historical interventions and does not claim that their behavioral shifts are a reproducible consequence of correctly implemented RL. The replication implementation passes temperature explicitly, uses the corrected extraction reward, and saves every rollout, extraction, reward, gradient checkpoint, and optimizer state. Its results remain separate from historical evaluation.

\begin{table*}[t]
\centering\small\begin{tabular}{lrlrr}
\toprule Condition & Eligible & Units & Amount: median [range] & Timing: median [range] \\
\midrule
Math & 4/14 & bits & +0.010 [-0.002, +0.043] & -0.014 [-0.030, +0.004] \\
Code & 14/14 & bits & +0.001 [-0.242, +0.130] & -0.060 [-0.116, -0.002] \\
GPQA & 14/14 & bits & +0.004 [-0.014, +0.036] & -0.013 [-0.047, +0.002] \\
Planning & 14/14 & bits & +0.020 [-0.014, +0.056] & -0.003 [-0.027, +0.024] \\
Moral & 14/14 & $\Delta R^2$ & -0.004 [-0.021, +0.045] & -0.010 [-0.035, +0.022] \\
Idea & 14/14 & $\Delta R^2$ & -0.008 [-0.050, +0.035] & -0.013 [-0.102, +0.014] \\
Safety & 6/14 & $\Delta R^2$ & -0.015 [-0.059, +0.034] & -0.079 [-0.324, -0.007] \\
Security & 14/14 & bits & +0.009 [-0.016, +0.043] & -0.013 [-0.052, +0.011] \\
\bottomrule
\end{tabular}

\caption{Eligible outcome-prediction tests and distribution of held-out increments. Ranges describe estimates across eligible model--family cells, not confidence intervals. The companion outcome table gives each estimate, 95\% interval, covariates, eligibility status, and adjusted $q$ value. No amount or timing test meets both the interval and FDR criteria.}
\label{tab:outcomes}
\end{table*}

\subsection{Corrected-reward bias learning}
\label{app:clean-rl}
The clean replication trains a single layer-16 residual bias in Qwen3.5-9B, starting from either zero or the frozen A2 vector, with three training seeds per initialization. Each run uses the same 301 training problems for 60 fixed update steps: 16 problems per step and eight sampled responses per problem, capped at 8,192 tokens. Adam uses learning rate 0.002. Sampling temperature is 0.6, with top-$p$ 0.95 and top-$k$ 20. The reward uses reference-blind Gemma4-31B-it answer extraction, grounding, and mathematical equivalence; unfinished generations receive zero. Within each problem, the mean reward is subtracted to form advantages. The update sums advantage-weighted generated-token log probabilities, normalized by batch size and the token cap. Temperature is passed explicitly; the likelihood uses the untruncated softmax even though sampling applies top-$k$/top-$p$ truncation. Thus this replication repairs the documented temperature and grading defects without claiming an exact on-policy likelihood for the truncated sampler.

Every run retains all 7,680 rollout--reward pairs, optimizer state, source hashes, and 61 bias vectors including initialization. The final checkpoint is step 60 for every seed, with no evaluation-based selection. Steps with zero advantages preserve the vector and remain in the fixed schedule. Independent audits verify complete seeded prompt/sample schedules, token bounds, finite vectors, and optimizer/checkpoint consistency.

Evaluation uses the same 199 historical test problems, two new sampled responses per problem, and an 8,192-token cap for all six learned biases and fresh unsteered/A2 controls. These problems are disjoint from the 301 training IDs but were used in earlier intervention-design comparisons; they are not a fresh project holdout. Accuracy counts unfinished attempts as incorrect. Paired intervals use 2,000 problem bootstraps after averaging the two samples. They are pointwise and conditional on each trained vector, not intervals over the population of training runs. All three seeds per initialization are reported separately.


All eight evaluation arms finish with the same 398 task--sample pairs (3,184 attempts total). All 2,388 learned-vector grades reproduce exactly from saved extractions, and all 26 paired accuracy/completion estimates and intervals reproduce independently. Across the eight arms, 18 extractions use fallback handling; companion tables bound each point estimate under arbitrary rescoring of those rows.

\begin{figure*}[t]
\centering\includegraphics[width=.98\textwidth]{figures/figA11_clean_rl.pdf}
\caption{\textbf{Corrected-reward learning improves natural completion; initialization matters.} Qwen3.5-9B, fixed step 60, 8,192 tokens, and all three training seeds per initialization. Paired differences from fresh unsteered (A) and fixed-A2 (B) controls use the same 199 historical evaluation problems and two samples per problem. Whiskers are pointwise 95\% problem-bootstrap intervals conditional on each learned vector. They do not measure training-population uncertainty. These evaluations do not include answer elicitation or annotation of the newly learned vectors.}
\label{fig:clean-rl}
\end{figure*}

Unsteered accuracy is 43.5\%, compared with 72.4\% for fixed A2. Zero-start RL attains 70.4\%, 68.8\%, and 71.6\% for seeds zero, one, and two: gains of 26.9, 25.4, and 28.1 pp over unsteered, with all three pointwise intervals above zero. None exceeds fixed A2; their accuracy-difference intervals include zero. A2-initialized RL attains 79.9\%, 81.2\%, and 77.4\%, exceeding fixed A2 by 7.5 pp [4.8, 10.3], 8.8 pp [5.8, 11.8], and 5.0 pp [2.5, 7.8], respectively. All three A2-initialized point-estimate gains remain positive even if every fallback grade in both comparison arms is reassigned adversely.

Completion follows the accuracy changes closely: it is 43.5\% unsteered, 73.4\% for fixed A2, 69.3--71.9\% after zero-start learning, and 78.9--83.7\% after A2-initialized learning. Thus the clean replication establishes improved natural-completion performance under this reward and budget, including further gains from learning around an atlas-derived initialization. It does not establish that the atlas supplies an optimal initialization, that learned gains survive common-budget elicitation, or that the new vectors reproduce the historical annotation shifts. The historical behavioral comparison below remains separate.

\subsection{Cross-family annotation robustness}
\label{app:multijudge}
The conversational and cognitive definitions come from \citet{kim2026society,gandhi2025habits}. \citet{kim2026society} report human agreement for both annotation families. Our indexed-sentence adaptation, judges, and task distribution differ from that validation. The robustness study compares Gemma4-31B-it, Qwen3.5-27B, and Llama-3.1-8B-Instruct, using unchanged templates and identical text. The third judge has smaller capacity; judge disagreement is reported rather than resolved by declaring a majority to be ground truth.

The frozen sample contains 1,381 traces: 805 atlas attempts, 256 paired intervention traces, and 320 matched-demand traces. Each atlas model--domain cell contributes 12 randomly sampled attempts plus an independent sample of up to four unfinished attempts; known union inclusion probabilities support weighted rates. Interventions use 16 common problems per target and one decoding seed. Matched demands use eight common problems per model and domain, including all four conditions. Text is restricted to the first 2,048 Qwen3.5-9B tokenizer units for every generator, so all judges see exactly the same prefix. Empty prefixes remain in coverage accounting. Generator names, treatment labels, original annotations, and outcomes are withheld from judges, although the text may reveal their source.

The frozen analysis reports per-label agreement, positive agreement, negative agreement, and Cohen's $\kappa$, with trace-cluster uncertainty and domain-specific disagreement. Positive agreement is undefined, not perfect, when a label is absent. Parse failures remain missing. Per-judge and common-valid-segment profiles test the sensitivity of domain contrasts, intervention shifts, and matched demands. This is a fixed-prefix robustness analysis; correlated judge errors remain possible. Exact prompts are included as companion text files. The original Llama pilot failed the 95\% cognitive parsing gate (92.0\%). Constraining JSON serialization to the requested sentence count/order and bounded whitespace raised valid cognitive parsing to 95.2\% and conversational parsing to 96.9\% on the same 69 nonempty traces. Prompts and label definitions are unchanged; original and repaired outputs are retained separately. On common-valid segments, label agreement is 97.5--99.7\%, with trace-mean rate changes below 0.6 pp. This comparison documents instrument sensitivity, not semantic validity.

All three judges cover the same 1,370 nonempty traces and 105,302 segments; 11 empty prefixes remain in the allocated denominator. There are no duplicated segments or context truncations. Conversational/cognitive parsing is 100/100\% for Gemma4-31B-it, 99.98/100\% for Qwen3.5-27B, and 97.91/95.86\% for Llama-3.1-8B-Instruct. Parsing measures usable output, not label correctness.

On atlas prefixes, Gemma--Qwen positive agreement ranges from 0.460 (reconciliation) to 0.790 (question--answering), with $\kappa$ from 0.447 to 0.775. Comparisons involving Llama have positive agreement of 0.013--0.347 and $\kappa$ of 0.009--0.324. In particular, question--answering and backward chaining have little agreement with Llama despite its high parsing coverage. The results do not support interchangeable sentence-level annotators.

The four selected domain contrasts nevertheless retain their average direction under every judge, both on judge-valid and common-valid segments. Using common-valid segments, idea generation exceeds math in perspective shifts by 19.2, 17.3, and 13.4 pp for Gemma, Qwen, and Llama; backtracking differences are 10.6, 8.6, and 2.5 pp. Planning exceeds moral reasoning in verification by 16.5, 13.4, and 6.3 pp, and math exceeds planning in subgoals by 9.3, 6.9, and 11.3 pp. Gemma and Qwen retain each direction in all seven generator models; Llama retains backtracking and verification in six of seven and the other two in all seven. These are descriptive checks of selected contrasts, not new significance tests or a reproduction of the full-trace timing analysis. Aggregate contrast stability can coexist with substantial event-level disagreement.

\begin{figure*}[t]
\centering\includegraphics[width=.98\textwidth]{figures/figA10_multijudge.pdf}
\caption{\textbf{Domain contrasts survive more readily than individual labels agree.} (A) Trace-balanced positive agreement on atlas prefixes, with pointwise 95\% trace-bootstrap intervals; the horizontal divider separates conversational and cognitive labels. Judge names are expanded in B. (B) Equal-generator mean differences after inverse-inclusion weighting within model/domain, restricted to segments valid for all three judges. Points are descriptive estimates, without uncertainty bars; counts give generator models with positive differences. All judges use identical fixed prefixes. No judge or majority consensus is treated as ground truth.}
\label{fig:multijudge}
\end{figure*}

Companion tables report all pairwise agreement metrics by domain and source, coverage by generator, both validity scopes, and paired intervention and matched-demand shifts. Independent checks reproduce all reported agreement point estimates, weighted rates, contrast means, and paired-shift means from the saved labels and profiles. Bootstrap intervals describe sampling uncertainty conditional on each annotation instrument; they do not quantify semantic validity.

\subsection{Matched-total-budget answer elicitation}
\label{app:elicitation}
The historical forced-answer probe retains the entire own answer once the thinking block closes; otherwise it adds a cue and up to 48 answer tokens. Thus its horizontal axis is a reasoning cutoff, not a common total-generation budget. The revised control reserves both cue tokens and the 48-token allowance inside the total budget, and retains an own answer only when the whole response fits. Frozen prefixes from baseline, domain vectors, and controls are answered by the same unsteered Qwen3.5-9B. Math includes A2, both saved RL vectors, matched SEAL, random, and concise controls. GPQA/planning include the math, GPQA, and idea contrasts and random control. This is retrospective protocol testing on the historical evaluation items.

All 27,920 allocated trace--budget rows are present: 15,920 math and 6,000 each for GPQA and planning. Each arm has the same task--seed--budget grid as its baseline; all rows satisfy the token bound. There are 199 math problems and 200 per other target, each with two decoding seeds. Intervals use 2,000 paired problem bootstraps after seed averaging and are pointwise, without multiplicity correction. The total budget counts generated prefix, cue, and answer tokens, excluding the common input prompt; it does not equate wall-clock time or prefill computation. Completed responses retain corrected extraction grades when available; forced answers use the frozen direct parser. Grading sources and all budget-specific estimates accompany the paper.

\begin{figure*}[t]
\centering\includegraphics[width=.95\textwidth]{figures/figA8_matched_elicitation.pdf}
\caption{\textbf{Large short-budget transfer gains do not survive common-budget answer elicitation.} Qwen3.5-9B, paired accuracy differences with pointwise 95\% problem-bootstrap intervals. Every displayed source--target pair is shown at both total budgets; each panel lists absolute baseline accuracy in the same order. Prefixes come from the frozen steered or baseline traces; any forced continuation uses the same unsteered model. All arms, including the math-specific controls, are reported in the companion tables.}
\label{fig:elicitation}
\end{figure*}

The math contrast gives GPQA accuracy of 69.5\% versus 76.5\% unsteered at 12,288 tokens, and planning accuracy of 92.3\% versus 91.0\%. At 32,768 tokens the corresponding differences are $-8.3$ pp [$-12.3$, $-4.5$] and $+0.5$ pp [$-2.8$, $3.8$]. The GPQA and idea contrasts remain adverse across all three targets at the two displayed budgets (Figure~\ref{fig:elicitation}). These results narrow the transfer claim to performance under the specified completion and elicitation procedures, rather than a general improvement from importing another domain's reasoning.

The math controls further distinguish fixed interventions. At 2,048 tokens, the saved zero-start RL vector improves accuracy by 5.8 pp [2.5, 9.3], versus $-0.5$ pp [$-4.3$, $3.3$] for A2, $+0.5$ pp [$-2.8$, $4.0$] for matched SEAL, and $+3.3$ pp [0.5, 6.5] for the concise prompt. At 32,768 tokens the zero-start difference is 1.3 pp [$-0.5$, $3.3$], and A2 is $-1.8$ pp [$-4.0$, $0.5$]. These are evaluations of historically selected vectors on reused items, not results of the corrected-reward replication.

Answer availability here means a nonempty extracted string, not semantic validity. In GPQA, 24 forced rows (20 lower-case \texttt{c} and four \texttt{f}) refer to labels embedded within questions; the frozen direct grader scores them zero. They remain in the denominator. Even changing all such rows in an arm--budget cell to correct could change a displayed GPQA contrast by at most 0.75 pp. We retain the frozen grades and report this sensitivity bound; the adverse math-to-GPQA conclusion is unaffected.

\subsection{Frozen fresh-item confirmation}
\label{app:confirmation}
A separate confirmation fixes A2 and the moral-associated contrast at their existing layer-16 norm of 6.75 before evaluating new items. The latter is selected for its strong retrospective math fidelity; neither is claimed to be a prospectively optimal choice. We draw 256 problems from the remaining MATH test set after excluding all 500 MATH-500 problems and screening against 740 distinct math prompt texts in the saved project task and trace inventory. Source revision, task hashes, exclusion rules, vector hashes, and decoding settings are frozen in the companion protocol. This is freshness relative to inventoried project use, not a model-pretraining decontamination claim.

Each arm and an unsteered baseline receive two sampled completions per problem at 8,192 tokens. The primary comparisons are the two natural-completion accuracy differences, with problem-level pairing after averaging samples and Holm correction. Common-total-budget elicitation, answer availability, token use, and literal repetition are secondary endpoints. A fixed 64-problem subset, selected before generation, supplies paired behavioral shifts and a fidelity check for the moral contrast. These analyses use no outcome-based sample or dose reselection. All 1,536 allocated attempts are present, with exact task--sample pairing across arms. Raw output hashes, frozen vectors, token accounting, own-answer grades, and all forced-answer parses and grades reproduce the saved outcomes. Intervals use 10,000 paired problem bootstraps; primary two-sided sign-flip tests use 10,000 draws and Holm correction over the two contrasts.

\begin{figure*}[t]
\centering\includegraphics[width=.95\textwidth]{figures/figA9_fresh_confirmation.pdf}
\caption{\textbf{Fresh items preserve completion gains, not an elicited-answer advantage.} Fixed A2 and moral-domain vectors on Qwen3.5-9B; two samples per problem. (A) Paired accuracy differences; both natural-completion contrasts pass the prespecified Holm correction, whereas elicited contrasts are secondary. (B) Future-operation readout targets reconstructed from annotations in the fixed 64-task subset. Whiskers are pointwise 95\% problem-bootstrap intervals. All 16 targets, including next-four rates, appear in the companion table.}
\label{fig:confirmation}
\end{figure*}

Natural-completion accuracy rises from 49.8\% unsteered to 82.6\% with A2 ($+32.8$ pp [27.9, 37.9], Holm $p<.001$) and 58.4\% with the moral contrast ($+8.6$ pp [3.3, 13.9], Holm $p=.002$). Completion rates are 50.4\%, 84.6\%, and 59.8\%, respectively. A2 reduces mean generated length by 1,997 tokens [1,765, 2,227]; the moral contrast changes it by $-91$ [$-303$, $129$]. Literal repetition falls from 7.4\% unsteered to 0.4\% and 2.3\%.

At the same total budget, answer elicitation yields accuracies of 95.1\%, 95.7\%, and 91.8\%. A2's difference is $+0.6$ pp [$-1.6$, $2.7$]; the moral contrast's is $-3.3$ pp [$-6.1$, $-0.6$]. These secondary intervals are pointwise. The baseline is near ceiling under this procedure, which limits detectable improvements; A2's interval does not establish equivalence. All 539 forced continuations yield a parseable answer. Four natural extractions use fallback handling (one baseline, one A2, two moral); flipping all affected grades could change a contrast by at most 0.6 pp. Frozen grades are retained. Elicitation uses the same unsteered reader and answer cue across arms, reserving cue and up to 48 answer tokens inside the 8,192-token cap. It changes the procedure rather than identifying the causal contribution of termination.

The behavioral audit retains all 192 traces (64 problems per arm, sample zero) and 55,588 segments. Both annotation families parse on every segment, with no context truncation or missing trace profiles; this is coverage evidence, not independent semantic validation. All 32 target-level mean shifts and their 10,000-replicate paired-bootstrap intervals reproduce from the exported profiles.

The moral contrast's achieved 16-target shift aligns with its intended direction (cosine 0.947 [0.847, 0.957]), with projection coefficient 0.450 [0.332, 0.570]. Thus it induces an aligned but smaller change, not exact source-profile reproduction. Its future-verification rate falls by 4.5 pp [2.0, 7.1], while future-subgoal rate rises by 3.7 pp [2.0, 5.4]. A2 raises future-subgoal rate by 6.8 pp [4.8, 8.8], with decreases in questioning, conflict, and reconciliation (Figure~\ref{fig:confirmation}B). These are average readout-target shifts rather than whole-trace prevalence changes. The moral vector therefore supplies a case where directional fidelity and improved natural completion coexist with lower elicited-answer accuracy. This does not identify behavioral mediation or validate a general intervention-selection rule.

\section{Outcome-Prediction Specification}
\label{app:outcome-spec}

Seven models, eight benchmark conditions, and two annotation families define 112 candidate cells. Cells require at least 40 scored traces; binary outcomes require at least 20 observations per class and continuous outcomes at least three distinct scores. Ten math cells lack binary-class support and eight safety cells lack outcome variation, leaving 94 eligible tests. Binary endpoints are correctness for math, code, GPQA, planning, and security; moral, idea, and safety retain their continuous rubric endpoints, oriented so larger is favorable under that rubric.

Each baseline includes log generated-token count, log valid-segment count, and available prompt-difficulty strata. Adding four operation rates tests amount; adding eight timing contrasts after those rates tests timing. Five frozen prompt folds and training-only imputation/scaling are used. Binary models are ridge logistic regressions; continuous models are ridge linear regressions. The penalty is one on slopes and zero on intercepts. Binary increments are changes in held-out log score in bits; continuous increments are squared-error reduction divided by outcome variance, equivalent to $\Delta R^2$. One thousand paired prompt bootstraps are stratified by class or outcome quantile. One-sided sign-flip tests receive Benjamini--Hochberg correction separately for the amount and timing families. Support requires a positive 95\% interval lower bound and $q<.05$. The estimates and their uncertainty distinguish this restricted null from a claim that profiles contain no success information.

\setcounter{dbltopnumber}{3}
\renewcommand{\dbltopfraction}{.95}
\renewcommand{\textfraction}{.03}
\section{Intervention Fidelity and Sensitivity Checks}
\label{app:fidelity}

We reconstruct the readout's future and next-four targets from valid annotated segments, using up to 80 segment endpoints inside the 8,192-token window. Achieved changes average paired trace-level means. We separately retain the intended contrast $\delta_s=\mu_s-\bar\mu$, fixed-map prediction $(v_s/\sigma)^\top W$, and observed change. Their cosine and projection onto $\delta_s$ measure directional alignment and magnitude; neither establishes mediation or successful reproduction of a source profile. Pairwise intervals resample prompts 2,000 times. Twenty-five source--target cells and their dose variants are exploratory, with shared baselines.

\begin{figure*}[t]
\centering\includegraphics[width=.90\textwidth]{figures/figA5_behavior_sensitivity.pdf}
\caption{\textbf{Average-shift similarity depends on window and magnitude.} Direction and size relative to A2 on 199 paired math items, using maximum absolute token windows. Bands are 95\% joint prompt-bootstrap intervals. Shorter completed responses contribute their observed text. Matched SEAL resembles A2 as closely as the saved zero-start RL vector at the longest window.}
\label{fig:behavior-sensitivity}
\end{figure*}
The joint bootstrap preserves the shared baseline in both intervention contrasts. A disjoint-item-half sensitivity uses different baseline observations for the two means: the median 8,192-token A2--RL cosine is 0.916 across random splits. This reduces literal shared-baseline dependence but does not remove differences in task composition or duration. A separate sensitivity conditions on all three responses reaching the window; at 2,048 tokens, 129 items remain and the cosine is 0.067 [$-0.410$, $0.571$]. This selected population is not a causal duration adjustment. The contrast with the full-population cosine shows that observed directional similarity is sensitive to population and window. Because magnitudes and coordinates also differ, these comparisons do not establish an explanatory insufficiency of the full representation.

\begin{figure*}[t]
\centering\includegraphics[width=.90\textwidth]{figures/figA6_demand_operations.pdf}
\caption{\textbf{All matched-demand conditions at the operation level.} Mean paired rate changes from free-response solving, in percentage points. Columns use Q, questioning; P, perspective shift; C, conflict; R, reconciliation; V, verification; B, backtracking; S, subgoals; K, backward chaining. Each model uses the same math/GPQA problems across answer-option, generate, and evaluate conditions; available pairs are reported in companion tables. Adding answer options changes available information as well as response format.}
\label{fig:demand-operations}
\end{figure*}
Generating approaches is the main demand contrast because both solve and generate have repeated-generation references. The evaluation condition is reported in Figure~\ref{fig:demand-operations}, without claiming an equivalent repeat-noise adjustment. The companion sensitivity table reports all contrasts at smoothing values $.001$, $.005$, and $.01$, together with mean absolute rate differences. These comparisons retain the same items and clarify which operations change; they do not isolate semantics from instruction-following conventions.

\section{Historical Fixed-Vector Comparisons}
\label{app:historical-vectors}
This secondary analysis compares a selected within-math archetype (A2), matched SEAL and random controls, and two saved biases from historical reward-learning runs, initialized from zero or A2. Only the layer-16 bias was optimized, using 301 math items disjoint from the 199 evaluation items, with one training seed per initialization. Evaluation uses corrected grades. The source audit identifies both an earlier reward parser and a temperature-variable defect; no per-run source hash or rollout log establishes their exact effect on training. We therefore analyze these as fixed interventions, not evidence for a reproducible tendency of correctly implemented RL (Appendix~\ref{app:audit}).

\begin{figure*}[t]
\centering\includegraphics[width=.95\textwidth,height=.47\textheight,keepaspectratio]{figures/fig05_steering_and_rl.pdf}
\caption{\textbf{Interventions differ in behavioral magnitude and performance despite similar shift directions.} Qwen3.5-9B, including saved historical RL vectors. (A) Paired first-8,192-token shifts: Q, questioning; P, perspective shift; C, conflict; R, reconciliation; V, verification; B, backtracking; S, subgoals; K, backward chaining. (B) Similarity to A2 in residual and mean behavior-change space; behavioral whiskers jointly bootstrap prompts. (C) Natural-completion accuracy, including matched SEAL and random controls. (D) Historical forced-answer differences with 95\% task-paired intervals; random is included, SEAL was not evaluated under this protocol. The RL training audit and answer-elicitation compute limit interpretation (Appendix~\ref{app:audit}).}
\label{fig:rl}
\end{figure*}
The saved zero-start RL vector has residual cosine approximately $-0.02$ with A2, but cosine 0.938 [0.839, 0.974] between their mean eight-rate shifts. The RL shift has only 0.412 times A2's magnitude. Both increase subgoals and reduce question--answer turns, conflict, and reconciliation. The matched SEAL direction is similarly aligned with A2 (0.934 [0.877, 0.968]), so the resemblance is not distinctive to reward learning. At a 1,024-token window, the A2--RL cosine is 0.740 [0.231, 0.911]. Joint prompt bootstraps, shorter windows, and disjoint-baseline sensitivities are reported in Appendix~\ref{app:fidelity}; these comparisons describe aggregate annotations rather than mechanisms. Restricting to the 129 items where baseline, A2, and the zero-start vector all reach 2,048 tokens gives cosine 0.067 [$-0.410$, $0.571$]. This post-treatment selection is not a causal duration adjustment, but shows substantial sensitivity to the observation window and population. Similar directions do not establish equivalent achieved profiles: magnitudes and individual coordinates differ.

At 8,192 tokens, natural-completion accuracy is 43.7\% unsteered, 73.6\% with A2, 67.3\% with zero-start RL, and 81.9\% with A2-initialized RL. At 32,768, the corresponding values are 86.4\%, 93.5\%, 97.2\%, and 94.5\%. Mean token use falls under all three interventions; these operating points do not establish a budget-independent best arm.

The historical forced-answer probe gives the saved zero-start RL vector a 7.3 pp [3.3, 11.6] advantage at a 1,024-token reasoning cutoff, where A2 does not show a comparable gain. This probe supplies a cue and additional decoding; it retains an entire final answer once thinking has closed, even if that answer extends beyond the cutoff. It therefore compares answer-elicitation protocols, not accuracy at identical total computation, and does not resolve the completion explanation for the domain-vector matrix. Appendix~\ref{app:elicitation} reports the completed common-total-budget control, including matched SEAL and concise prompting.



